% =====================================================================
%  Relativistic extension of Chapter II, Sections 13--14
%  (Variational Principles & Thermodynamic Significance)
%
%  Based on Chandrasekhar, "Hydrodynamic and Hydromagnetic Stability"
%  Original: §13 -- The variational principles
%            §14 -- The thermodynamic significance of the variational
%                   principle
%
%  Relativistic formulation following the BDNK (Bemfica--Disconzi--Noronha--Kovtun)
%  causal first-order dissipative framework; metric signature (-,+,+,+); c kept explicit.
%  References: Bemfica et al. (2018), Kovtun (2019).
% =====================================================================
% Relativistic Modification Preamble
% Shared macros for all relativistic chapter modifications
%
% NOTE: This file requires \usepackage{tcolorbox} to be loaded
% in the main preamble (../../preamble.tex) for the custom
% environments defined at the end of this file.
%
% Usage: % Relativistic Modification Preamble
% Shared macros for all relativistic chapter modifications
%
% NOTE: This file requires \usepackage{tcolorbox} to be loaded
% in the main preamble (../../preamble.tex) for the custom
% environments defined at the end of this file.
%
% Usage: % Relativistic Modification Preamble
% Shared macros for all relativistic chapter modifications
%
% NOTE: This file requires \usepackage{tcolorbox} to be loaded
% in the main preamble (../../preamble.tex) for the custom
% environments defined at the end of this file.
%
% Usage: \input{rel_preamble} from any relativistic chapter
% or from rel_main.tex

% =====================================================================
% 4-Vector notation
% ---------------------------------------------------------------------
% Macros for writing 4-vectors and projection operators in a
% consistent style throughout all relativistic chapters.
% =====================================================================
\newcommand{\fourv}[1]{#1^{\mu}}
\newcommand{\fvel}{u^{\mu}}
\newcommand{\fveldown}{u_{\mu}}
\newcommand{\proj}{\Delta^{\mu\nu}}
\newcommand{\projdown}{\Delta_{\mu\nu}}

% =====================================================================
% Tensors
% ---------------------------------------------------------------------
% Energy-momentum tensor, Faraday tensor and its dual, metric tensors.
% These appear pervasively in the relativistic formulation of both
% hydrodynamic and hydromagnetic stability problems.
% =====================================================================
\newcommand{\emt}{T^{\mu\nu}}
\newcommand{\emtdown}{T_{\mu\nu}}
\newcommand{\farad}{F^{\mu\nu}}
\newcommand{\faraddown}{F_{\mu\nu}}
\newcommand{\dualF}{{}^{*}\!F^{\mu\nu}}
\newcommand{\metric}{g^{\mu\nu}}
\newcommand{\metricdown}{g_{\mu\nu}}
\newcommand{\mink}{\eta^{\mu\nu}}

% =====================================================================
% Covariant derivatives and Lie derivative
% ---------------------------------------------------------------------
% Covariant derivatives with different index labels for use in
% multi-index expressions without conflict.
% =====================================================================
\newcommand{\covd}{\nabla_{\mu}}
\newcommand{\covdu}{\nabla_{\nu}}
\newcommand{\covda}{\nabla_{\alpha}}
\newcommand{\Lie}{\mathcal{L}}

% =====================================================================
% Physical quantities
% ---------------------------------------------------------------------
% Core thermodynamic and kinematic quantities that appear in the
% relativistic equations of motion.  The enthalpy density
% w = epsilon + p plays the role that rho alone plays in the
% Newtonian theory.
% =====================================================================
\newcommand{\enthalpy}{w}           % w = epsilon + p
\newcommand{\edensity}{\varepsilon} % energy density
\newcommand{\rdensity}{\rho_0}     % rest-mass density
\newcommand{\Lf}{\gamma}            % Lorentz factor
\newcommand{\cs}{c_{\mathrm{s}}}    % sound speed
\newcommand{\vA}{v_{\mathrm{A}}}    % Alfven speed
\newcommand{\vf}{v_{\mathrm{f}}}    % fast magnetosonic speed
\newcommand{\vs}{v_{\mathrm{s}}}    % slow magnetosonic speed
\newcommand{\bfour}{b^{\mu}}        % magnetic 4-vector
\newcommand{\bsq}{b^2}             % b^mu b_mu (magnetic pressure term)

% =====================================================================
% Dissipation (Israel-Stewart causal theory)
% ---------------------------------------------------------------------
% The Israel-Stewart formalism replaces the acausal Navier-Stokes
% constitutive relations with hyperbolic relaxation equations.
% Each dissipative flux (bulk pressure Pi, shear tensor pi^{mu nu},
% heat flux q^mu) has its own relaxation time (tau_Pi, tau_pi, tau_q).
% =====================================================================
\newcommand{\bulkP}{\Pi}            % bulk viscous pressure
\newcommand{\shearten}{\pi^{\mu\nu}} % shear stress tensor
\newcommand{\heatflux}{q^{\mu}}     % heat flux 4-vector
\newcommand{\taupi}{\tau_{\pi}}     % shear relaxation time
\newcommand{\tauq}{\tau_{q}}        % heat relaxation time
\newcommand{\tauPi}{\tau_{\Pi}}     % bulk relaxation time
\newcommand{\bulkvisc}{\zeta}       % bulk viscosity coefficient
\newcommand{\shearvisc}{\eta_{\mathrm{s}}} % shear viscosity coefficient
\newcommand{\thermcond}{\kappa}     % thermal conductivity

% =====================================================================
% Dimensionless numbers (relativistic generalisations)
% ---------------------------------------------------------------------
% Relativistic analogues of the classical dimensionless parameters
% that govern stability thresholds.  Subscript "rel" distinguishes
% them from the Newtonian versions defined in the original chapters.
% =====================================================================
\newcommand{\Rarel}{\mathrm{Ra}_{\mathrm{rel}}}   % relativistic Rayleigh number
\newcommand{\Tarel}{\mathrm{Ta}_{\mathrm{rel}}}   % relativistic Taylor number
\newcommand{\Qrel}{Q_{\mathrm{rel}}}              % relativistic Chandrasekhar Q
\newcommand{\Mach}{\mathcal{M}}                    % Mach number

% =====================================================================
% Operators
% ---------------------------------------------------------------------
% The d'Alembertian (wave operator) replaces the Laplacian in the
% relativistic wave equations.  Two aliases are provided for
% author convenience.
% =====================================================================
\newcommand{\dAlembert}{\Box}       % d'Alembertian
\newcommand{\Dalembertian}{\Box}    % alias
\newcommand{\order}[1]{\mathcal{O}(#1)} % order-of-magnitude symbol

% =====================================================================
% Relativistic correction helpers
% ---------------------------------------------------------------------
% \relcorr{v^2/c^2} expands to [1 + O(v^2/c^2)], useful for
% writing correction factors.  \NRlimit produces an arrow
% labelled "v/c -> 0" for taking the non-relativistic limit.
% =====================================================================
\newcommand{\relcorr}[1]{\left[1 + \order{#1}\right]} % correction factor
\newcommand{\NRlimit}{\xrightarrow{v/c \to 0}}        % non-relativistic limit arrow

% =====================================================================
% Custom environments (require tcolorbox package)
% ---------------------------------------------------------------------
% causalitycheck  -- yellow box for verifying that perturbation
%                    equations preserve causality (signal speeds <= c).
% relcorrection   -- blue box highlighting where and how the
%                    relativistic result departs from the Newtonian one.
% =====================================================================
\newenvironment{causalitycheck}{%
  \begin{tcolorbox}[colback=yellow!5, colframe=yellow!50!black, title=Causality Verification]
}{%
  \end{tcolorbox}
}

\newenvironment{relcorrection}{%
  \begin{tcolorbox}[colback=blue!5, colframe=blue!50!black, title=Relativistic Correction]
}{%
  \end{tcolorbox}
}
 from any relativistic chapter
% or from rel_main.tex

% =====================================================================
% 4-Vector notation
% ---------------------------------------------------------------------
% Macros for writing 4-vectors and projection operators in a
% consistent style throughout all relativistic chapters.
% =====================================================================
\newcommand{\fourv}[1]{#1^{\mu}}
\newcommand{\fvel}{u^{\mu}}
\newcommand{\fveldown}{u_{\mu}}
\newcommand{\proj}{\Delta^{\mu\nu}}
\newcommand{\projdown}{\Delta_{\mu\nu}}

% =====================================================================
% Tensors
% ---------------------------------------------------------------------
% Energy-momentum tensor, Faraday tensor and its dual, metric tensors.
% These appear pervasively in the relativistic formulation of both
% hydrodynamic and hydromagnetic stability problems.
% =====================================================================
\newcommand{\emt}{T^{\mu\nu}}
\newcommand{\emtdown}{T_{\mu\nu}}
\newcommand{\farad}{F^{\mu\nu}}
\newcommand{\faraddown}{F_{\mu\nu}}
\newcommand{\dualF}{{}^{*}\!F^{\mu\nu}}
\newcommand{\metric}{g^{\mu\nu}}
\newcommand{\metricdown}{g_{\mu\nu}}
\newcommand{\mink}{\eta^{\mu\nu}}

% =====================================================================
% Covariant derivatives and Lie derivative
% ---------------------------------------------------------------------
% Covariant derivatives with different index labels for use in
% multi-index expressions without conflict.
% =====================================================================
\newcommand{\covd}{\nabla_{\mu}}
\newcommand{\covdu}{\nabla_{\nu}}
\newcommand{\covda}{\nabla_{\alpha}}
\newcommand{\Lie}{\mathcal{L}}

% =====================================================================
% Physical quantities
% ---------------------------------------------------------------------
% Core thermodynamic and kinematic quantities that appear in the
% relativistic equations of motion.  The enthalpy density
% w = epsilon + p plays the role that rho alone plays in the
% Newtonian theory.
% =====================================================================
\newcommand{\enthalpy}{w}           % w = epsilon + p
\newcommand{\edensity}{\varepsilon} % energy density
\newcommand{\rdensity}{\rho_0}     % rest-mass density
\newcommand{\Lf}{\gamma}            % Lorentz factor
\newcommand{\cs}{c_{\mathrm{s}}}    % sound speed
\newcommand{\vA}{v_{\mathrm{A}}}    % Alfven speed
\newcommand{\vf}{v_{\mathrm{f}}}    % fast magnetosonic speed
\newcommand{\vs}{v_{\mathrm{s}}}    % slow magnetosonic speed
\newcommand{\bfour}{b^{\mu}}        % magnetic 4-vector
\newcommand{\bsq}{b^2}             % b^mu b_mu (magnetic pressure term)

% =====================================================================
% Dissipation (Israel-Stewart causal theory)
% ---------------------------------------------------------------------
% The Israel-Stewart formalism replaces the acausal Navier-Stokes
% constitutive relations with hyperbolic relaxation equations.
% Each dissipative flux (bulk pressure Pi, shear tensor pi^{mu nu},
% heat flux q^mu) has its own relaxation time (tau_Pi, tau_pi, tau_q).
% =====================================================================
\newcommand{\bulkP}{\Pi}            % bulk viscous pressure
\newcommand{\shearten}{\pi^{\mu\nu}} % shear stress tensor
\newcommand{\heatflux}{q^{\mu}}     % heat flux 4-vector
\newcommand{\taupi}{\tau_{\pi}}     % shear relaxation time
\newcommand{\tauq}{\tau_{q}}        % heat relaxation time
\newcommand{\tauPi}{\tau_{\Pi}}     % bulk relaxation time
\newcommand{\bulkvisc}{\zeta}       % bulk viscosity coefficient
\newcommand{\shearvisc}{\eta_{\mathrm{s}}} % shear viscosity coefficient
\newcommand{\thermcond}{\kappa}     % thermal conductivity

% =====================================================================
% Dimensionless numbers (relativistic generalisations)
% ---------------------------------------------------------------------
% Relativistic analogues of the classical dimensionless parameters
% that govern stability thresholds.  Subscript "rel" distinguishes
% them from the Newtonian versions defined in the original chapters.
% =====================================================================
\newcommand{\Rarel}{\mathrm{Ra}_{\mathrm{rel}}}   % relativistic Rayleigh number
\newcommand{\Tarel}{\mathrm{Ta}_{\mathrm{rel}}}   % relativistic Taylor number
\newcommand{\Qrel}{Q_{\mathrm{rel}}}              % relativistic Chandrasekhar Q
\newcommand{\Mach}{\mathcal{M}}                    % Mach number

% =====================================================================
% Operators
% ---------------------------------------------------------------------
% The d'Alembertian (wave operator) replaces the Laplacian in the
% relativistic wave equations.  Two aliases are provided for
% author convenience.
% =====================================================================
\newcommand{\dAlembert}{\Box}       % d'Alembertian
\newcommand{\Dalembertian}{\Box}    % alias
\newcommand{\order}[1]{\mathcal{O}(#1)} % order-of-magnitude symbol

% =====================================================================
% Relativistic correction helpers
% ---------------------------------------------------------------------
% \relcorr{v^2/c^2} expands to [1 + O(v^2/c^2)], useful for
% writing correction factors.  \NRlimit produces an arrow
% labelled "v/c -> 0" for taking the non-relativistic limit.
% =====================================================================
\newcommand{\relcorr}[1]{\left[1 + \order{#1}\right]} % correction factor
\newcommand{\NRlimit}{\xrightarrow{v/c \to 0}}        % non-relativistic limit arrow

% =====================================================================
% Custom environments (require tcolorbox package)
% ---------------------------------------------------------------------
% causalitycheck  -- yellow box for verifying that perturbation
%                    equations preserve causality (signal speeds <= c).
% relcorrection   -- blue box highlighting where and how the
%                    relativistic result departs from the Newtonian one.
% =====================================================================
\newenvironment{causalitycheck}{%
  \begin{tcolorbox}[colback=yellow!5, colframe=yellow!50!black, title=Causality Verification]
}{%
  \end{tcolorbox}
}

\newenvironment{relcorrection}{%
  \begin{tcolorbox}[colback=blue!5, colframe=blue!50!black, title=Relativistic Correction]
}{%
  \end{tcolorbox}
}
 from any relativistic chapter
% or from rel_main.tex

% =====================================================================
% 4-Vector notation
% ---------------------------------------------------------------------
% Macros for writing 4-vectors and projection operators in a
% consistent style throughout all relativistic chapters.
% =====================================================================
\newcommand{\fourv}[1]{#1^{\mu}}
\newcommand{\fvel}{u^{\mu}}
\newcommand{\fveldown}{u_{\mu}}
\newcommand{\proj}{\Delta^{\mu\nu}}
\newcommand{\projdown}{\Delta_{\mu\nu}}

% =====================================================================
% Tensors
% ---------------------------------------------------------------------
% Energy-momentum tensor, Faraday tensor and its dual, metric tensors.
% These appear pervasively in the relativistic formulation of both
% hydrodynamic and hydromagnetic stability problems.
% =====================================================================
\newcommand{\emt}{T^{\mu\nu}}
\newcommand{\emtdown}{T_{\mu\nu}}
\newcommand{\farad}{F^{\mu\nu}}
\newcommand{\faraddown}{F_{\mu\nu}}
\newcommand{\dualF}{{}^{*}\!F^{\mu\nu}}
\newcommand{\metric}{g^{\mu\nu}}
\newcommand{\metricdown}{g_{\mu\nu}}
\newcommand{\mink}{\eta^{\mu\nu}}

% =====================================================================
% Covariant derivatives and Lie derivative
% ---------------------------------------------------------------------
% Covariant derivatives with different index labels for use in
% multi-index expressions without conflict.
% =====================================================================
\newcommand{\covd}{\nabla_{\mu}}
\newcommand{\covdu}{\nabla_{\nu}}
\newcommand{\covda}{\nabla_{\alpha}}
\newcommand{\Lie}{\mathcal{L}}

% =====================================================================
% Physical quantities
% ---------------------------------------------------------------------
% Core thermodynamic and kinematic quantities that appear in the
% relativistic equations of motion.  The enthalpy density
% w = epsilon + p plays the role that rho alone plays in the
% Newtonian theory.
% =====================================================================
\newcommand{\enthalpy}{w}           % w = epsilon + p
\newcommand{\edensity}{\varepsilon} % energy density
\newcommand{\rdensity}{\rho_0}     % rest-mass density
\newcommand{\Lf}{\gamma}            % Lorentz factor
\newcommand{\cs}{c_{\mathrm{s}}}    % sound speed
\newcommand{\vA}{v_{\mathrm{A}}}    % Alfven speed
\newcommand{\vf}{v_{\mathrm{f}}}    % fast magnetosonic speed
\newcommand{\vs}{v_{\mathrm{s}}}    % slow magnetosonic speed
\newcommand{\bfour}{b^{\mu}}        % magnetic 4-vector
\newcommand{\bsq}{b^2}             % b^mu b_mu (magnetic pressure term)

% =====================================================================
% Dissipation (Israel-Stewart causal theory)
% ---------------------------------------------------------------------
% The Israel-Stewart formalism replaces the acausal Navier-Stokes
% constitutive relations with hyperbolic relaxation equations.
% Each dissipative flux (bulk pressure Pi, shear tensor pi^{mu nu},
% heat flux q^mu) has its own relaxation time (tau_Pi, tau_pi, tau_q).
% =====================================================================
\newcommand{\bulkP}{\Pi}            % bulk viscous pressure
\newcommand{\shearten}{\pi^{\mu\nu}} % shear stress tensor
\newcommand{\heatflux}{q^{\mu}}     % heat flux 4-vector
\newcommand{\taupi}{\tau_{\pi}}     % shear relaxation time
\newcommand{\tauq}{\tau_{q}}        % heat relaxation time
\newcommand{\tauPi}{\tau_{\Pi}}     % bulk relaxation time
\newcommand{\bulkvisc}{\zeta}       % bulk viscosity coefficient
\newcommand{\shearvisc}{\eta_{\mathrm{s}}} % shear viscosity coefficient
\newcommand{\thermcond}{\kappa}     % thermal conductivity

% =====================================================================
% Dimensionless numbers (relativistic generalisations)
% ---------------------------------------------------------------------
% Relativistic analogues of the classical dimensionless parameters
% that govern stability thresholds.  Subscript "rel" distinguishes
% them from the Newtonian versions defined in the original chapters.
% =====================================================================
\newcommand{\Rarel}{\mathrm{Ra}_{\mathrm{rel}}}   % relativistic Rayleigh number
\newcommand{\Tarel}{\mathrm{Ta}_{\mathrm{rel}}}   % relativistic Taylor number
\newcommand{\Qrel}{Q_{\mathrm{rel}}}              % relativistic Chandrasekhar Q
\newcommand{\Mach}{\mathcal{M}}                    % Mach number

% =====================================================================
% Operators
% ---------------------------------------------------------------------
% The d'Alembertian (wave operator) replaces the Laplacian in the
% relativistic wave equations.  Two aliases are provided for
% author convenience.
% =====================================================================
\newcommand{\dAlembert}{\Box}       % d'Alembertian
\newcommand{\Dalembertian}{\Box}    % alias
\newcommand{\order}[1]{\mathcal{O}(#1)} % order-of-magnitude symbol

% =====================================================================
% Relativistic correction helpers
% ---------------------------------------------------------------------
% \relcorr{v^2/c^2} expands to [1 + O(v^2/c^2)], useful for
% writing correction factors.  \NRlimit produces an arrow
% labelled "v/c -> 0" for taking the non-relativistic limit.
% =====================================================================
\newcommand{\relcorr}[1]{\left[1 + \order{#1}\right]} % correction factor
\newcommand{\NRlimit}{\xrightarrow{v/c \to 0}}        % non-relativistic limit arrow

% =====================================================================
% Custom environments (require tcolorbox package)
% ---------------------------------------------------------------------
% causalitycheck  -- yellow box for verifying that perturbation
%                    equations preserve causality (signal speeds <= c).
% relcorrection   -- blue box highlighting where and how the
%                    relativistic result departs from the Newtonian one.
% =====================================================================
\newenvironment{causalitycheck}{%
  \begin{tcolorbox}[colback=yellow!5, colframe=yellow!50!black, title=Causality Verification]
}{%
  \end{tcolorbox}
}

\newenvironment{relcorrection}{%
  \begin{tcolorbox}[colback=blue!5, colframe=blue!50!black, title=Relativistic Correction]
}{%
  \end{tcolorbox}
}


\section{Relativistic variational principles}\label{sec:rel-13}

In the classical theory (\S~13), the critical Rayleigh number is
characterised as the stationary value of a ratio of positive-definite
integrals---the Rayleigh quotient.  We now derive the relativistic
analogues of these variational principles for a layer of fluid in a
static spacetime with a uniform gravitational field~$g$.

Throughout this section we work in Eulerian perturbation theory around
a relativistic hydrostatic equilibrium with energy density
$\edensity$, pressure $p$, and relativistic enthalpy density
$\enthalpy = \edensity + p$.  Length is measured in units of the layer
depth~$d$, and $D \equiv \mathrm{d}/\mathrm{d}z$ as before.  The
perturbation amplitude for the vertical component of the 4-velocity is
$\delta u^z \equiv W(z)\,f(x,y)\,e^{\sigma t}$, and the temperature
perturbation is $\Theta(z)\,f(x,y)\,e^{\sigma t}$, where $f$ satisfies
$(\partial_x^2+\partial_y^2)f = -a^2 f$.

\subsection{Relativistic characteristic value problem}\label{sec:rel-13-prelim}

The linearised relativistic perturbation equations for a thermally
stratified fluid, in the BDNK first-order causal framework
(Bemfica et al.\ 2018; Kovtun 2019) at marginal stability ($\sigma=0$), reduce
to (see \S\S~11--12 for the classical derivation, and the relativistic
linearisation in the companion sections)
\begin{equation}
\label{eq:rel-2-134}
(D^2-a^2)F_{\mathrm{rel}} = -\Rarel\;a^2\,\frac{\enthalpy}{c^2}\,W,
\end{equation}
where
\begin{equation}
\label{eq:rel-2-133}
F_{\mathrm{rel}} = (D^2-a^2)^2 W = (D^2-a^2)G_{\mathrm{rel}},
\qquad
G_{\mathrm{rel}} = (D^2-a^2)W,
\end{equation}
and the relativistic Rayleigh number is
\begin{equation}
\label{eq:rel-Ra-def}
\Rarel \;=\; \frac{\alpha\,g\,\beta\,d^4}{\kappa_T\,\nu}
\left(1 + \frac{\edensity + p}{\rdensity\,c^2}\right)
\;=\; \Ra\;\frac{\enthalpy}{\rdensity\,c^2}.
\end{equation}
Here $\alpha$ is the thermal expansion coefficient, $\beta$ the
adverse temperature gradient, $\kappa_T$ the thermal diffusivity, and
$\nu$ the kinematic viscosity.  The factor
$\enthalpy/(\rdensity c^2)$ is the key relativistic correction:
in the non-relativistic limit $\enthalpy \to \rdensity c^2$ and
$\Rarel \to \Ra$.

The boundary conditions are unchanged in form:
\begin{equation}
\label{eq:rel-2-135}
W = 0 \;\text{ and }\; F_{\mathrm{rel}} = 0
\;\text{ for }\; z = 0,\;1;\quad
\text{and \emph{either} } DW = 0 \;\text{ \emph{or} }\; D^2 W = 0,
\end{equation}
depending on whether the bounding surfaces are rigid or free.

\subsection{The first variational principle (relativistic)}%
\label{sec:rel-13a}

Let $\Rarel_j$ denote the characteristic values and let subscripts
$i,j$ label eigenfunctions belonging to different characteristic
values.  Multiplying the equation for $W_i$ by $F_{\mathrm{rel},j}$,
integrating, and performing the same sequence of integrations by parts
as in \S~13\textit{a}, we obtain the relativistic orthogonality
relation
\begin{equation}
\label{eq:rel-2-139}
\int_0^1 \frac{\enthalpy}{c^2}\,G_{\mathrm{rel},i}\,
G_{\mathrm{rel},j}\,\mathrm{d}z = 0
\qquad (i \neq j).
\end{equation}
The weight function $\enthalpy/c^2$ replaces the unit weight of the
classical theory; physically it reflects the relativistic inertia of
enthalpy.

Setting $i = j$ we find the \emph{relativistic Rayleigh quotient}:
\begin{equation}
\label{eq:rel-2-141}
\boxed{
\Rarel \;=\;
\frac{\displaystyle\int_0^1\bigl[(DF_{\mathrm{rel}})^2
  + a^2 F_{\mathrm{rel}}^2\bigr]\,\mathrm{d}z}
{a^2\;\displaystyle\int_0^1
  \frac{\enthalpy}{c^2}\,G_{\mathrm{rel}}^2\,\mathrm{d}z}
\;=\; \frac{\mathscr{I}_1^{\mathrm{rel}}}
           {\mathscr{I}_2^{\mathrm{rel}}}.
}
\end{equation}
Equation~\eqref{eq:rel-2-141} reduces to the classical
expression~(2-141) when $\enthalpy/c^2 \to \rdensity$ (i.e.\ when
internal and kinetic energies are negligible compared with rest-mass
energy).

\paragraph{Stationary property.}
Let $W$ be an admissible trial function satisfying~\eqref{eq:rel-2-135},
and let $\delta W$ be a variation compatible with those boundary
conditions.  Following the classical argument (equations (2-143)--(2-148)), we compute
\begin{equation}
\label{eq:rel-2-148}
\delta\Rarel = -\frac{2}{\mathscr{I}_2^{\mathrm{rel}}}
\int_0^1 \delta F_{\mathrm{rel}}
\left\{(D^2-a^2)F_{\mathrm{rel}}
  + \Rarel\,a^2\,\frac{\enthalpy}{c^2}\,W\right\}\mathrm{d}z.
\end{equation}
Hence $\delta\Rarel = 0$ if and only if~\eqref{eq:rel-2-134} holds,
establishing the stationary property.

\paragraph{Minimum property.}
Let the eigenfunctions $G_{\mathrm{rel},j}$ be normalised with respect
to the relativistic weight:
\begin{equation}
\label{eq:rel-2-150}
\int_0^1 \frac{\enthalpy}{c^2}\,
G_{\mathrm{rel},i}\,G_{\mathrm{rel},j}\,\mathrm{d}z
= \delta_{ij}.
\end{equation}
Expanding an arbitrary normalised $G_{\mathrm{rel}}$ in the
eigenbasis, $G_{\mathrm{rel}} = \sum_j A_j\,G_{\mathrm{rel},j}$,
and repeating the argument of equations (2-154)--(2-161), we obtain
\begin{equation}
\label{eq:rel-2-161}
\Rarel_{1}\;a^2 \;\leqslant\;
\int_0^1\bigl[(DF_{\mathrm{rel}})^2 + a^2 F_{\mathrm{rel}}^2\bigr]
\,\mathrm{d}z,
\end{equation}
with equality if and only if $A_j = 0$ for $j \geqslant 2$.
\emph{The lowest relativistic Rayleigh number is therefore a true
minimum of the relativistic Rayleigh quotient.}

\subsection{The second variational principle (relativistic)}%
\label{sec:rel-13b}

Starting instead from the characteristic value problem expressed in
terms of the temperature perturbation~$\Theta$, we let
$G'_{\mathrm{rel}} = (D^2 - a^2)\Theta$ and follow the procedure of
\S~13\textit{b}.  The relativistic second variational principle gives
\begin{equation}
\label{eq:rel-2-170}
\boxed{
\Rarel_j \;=\;
\frac{\displaystyle\int_0^1
  \bigl[(D^2-a^2)G'_{\mathrm{rel},j}\bigr]^2\,\mathrm{d}z}
{a^2\;\displaystyle\int_0^1
  \frac{\enthalpy}{c^2}\,
  \bigl[(D\Theta_j)^2 + a^2\Theta_j^2\bigr]\,\mathrm{d}z}.
}
\end{equation}
Again, the relativistic enthalpy weight $\enthalpy/c^2$ appears in
the denominator, and the same stationary and minimum properties hold.

\begin{relcorrection}
Both variational principles acquire a single structural modification
in the relativistic case: every integral that classically involves
the fluid density $\rho$ is replaced by an integral weighted with
the relativistic inertial mass density $\enthalpy/c^2 = (\edensity+p)/c^2$.
This replacement originates from the relativistic momentum equation
$\nabla_\nu T^{\mu\nu} = 0$, where the inertia of enthalpy---not
just rest mass---enters the equation of motion.
\end{relcorrection}

% =====================================================================
\section{Thermodynamic significance of the relativistic variational
principle}\label{sec:rel-14}

In \S~14, Chandrasekhar showed that the classical variational
principle has a direct thermodynamic interpretation: the Rayleigh
number is the ratio of the rate of energy release by buoyancy to the
rate of viscous dissipation.  We now establish the relativistic
counterpart of this result within the BDNK causal first-order
dissipative theory (Bemfica, Disconzi \& Noronha 2018; Kovtun 2019).

\subsection{Relativistic viscous dissipation rate}\label{sec:rel-14a}

At marginal stability ($\sigma = 0$) in the relativistic theory, the
$z$-component of the 4-vorticity vanishes by the same argument as in
the classical case (cf.\ equation~(2-171)).  The perturbation velocity
field therefore has the same kinematic structure as in
equations~(2-172)--(2-175), with $\rho$ replaced by $\enthalpy/c^2$.

The rate of viscous dissipation per unit horizontal area in the
relativistic theory is
\begin{equation}
\label{eq:rel-2-176}
\epsilon_v^{\mathrm{rel}} = -\frac{\enthalpy\,\nu}{c^2\,d^2}
\int_0^1 \bigl\langle
  w(D^2-a^2)w + u(D^2-a^2)u + v(D^2-a^2)v
\bigr\rangle\,\mathrm{d}z,
\end{equation}
where angle brackets denote the horizontal average and $\nu$ is the
kinematic shear viscosity.  The factor $\enthalpy/c^2$ replaces
the classical $\rho$ because the relativistic viscous stress
$\pi^{\mu\nu}$ acts on the inertia of enthalpy.

Following exactly the integration-by-parts sequence of
equations~(2-177)--(2-178), we obtain
\begin{equation}
\label{eq:rel-2-178}
\epsilon_v^{\mathrm{rel}}
= \frac{\enthalpy\,\nu}{4\,c^2\,d^2}
\int_0^1 G_{\mathrm{rel}}^2\,\mathrm{d}z
= \frac{\enthalpy\,\nu}{4\,c^2\,d^2}
\int_0^1 \bigl[(D^2-a^2)W\bigr]^2\,\mathrm{d}z.
\end{equation}

\subsection{Relativistic buoyancy energy release rate}\label{sec:rel-14b}

In the relativistic framework the buoyancy force per unit volume is
\begin{equation}
\label{eq:rel-buoy-force}
f_{\mathrm{buoy}}
= -\frac{\enthalpy}{c^2}\,g\,\alpha\,\theta,
\end{equation}
where $\theta$ is the Eulerian temperature perturbation and
$\alpha = -(\rdensity)^{-1}(\partial\rdensity/\partial T)_p$ the
thermal expansion coefficient.  The rate of energy release by buoyancy
in a unit column is therefore
\begin{equation}
\label{eq:rel-2-179}
\epsilon_g^{\mathrm{rel}}
= \frac{\enthalpy}{c^2}\,g\,\alpha
\int_0^1 \langle\theta\,w\rangle\,\mathrm{d}z.
\end{equation}
Using the linearised heat equation (the relativistic analogue of
equation~(2-76)) and integrating by parts gives
\begin{equation}
\label{eq:rel-2-181}
\epsilon_g^{\mathrm{rel}}
= \frac{\enthalpy\,g\,\alpha\,\kappa_T}
       {4\,c^2\,\beta\,d^2}
\int_0^1 \bigl[(D\Theta)^2 + a^2\Theta^2\bigr]\,\mathrm{d}z.
\end{equation}
Expressing $\Theta$ in terms of $F_{\mathrm{rel}}$ via the
relativistic analogue of equation~(2-182),
\begin{equation}
\label{eq:rel-2-182}
\Theta = \frac{\nu\,c^2}{g\,\alpha\,\enthalpy\,d^2}\,
(D^2-a^2)^2 W
= \frac{\nu\,c^2}{g\,\alpha\,\enthalpy\,d^2}\,F_{\mathrm{rel}},
\end{equation}
we obtain
\begin{equation}
\label{eq:rel-2-183}
\epsilon_g^{\mathrm{rel}}
= \frac{\nu^2\,\kappa_T\,c^2}
       {4\,g\,\alpha\,\beta\,\enthalpy\,d^6}
\int_0^1 \bigl[(DF_{\mathrm{rel}})^2
  + a^2 F_{\mathrm{rel}}^2\bigr]\,\mathrm{d}z.
\end{equation}

\subsection{Energy balance and the variational principle}\label{sec:rel-14c}

In a steady state the viscous dissipation and buoyancy release must
balance:
\begin{equation}
\label{eq:rel-2-184}
\epsilon_v^{\mathrm{rel}} = \epsilon_g^{\mathrm{rel}}.
\end{equation}
Substituting equations~\eqref{eq:rel-2-178}
and~\eqref{eq:rel-2-183} we find
\begin{equation}
\label{eq:rel-2-185}
\frac{\enthalpy\,\nu}{4\,c^2\,d^2}
\int_0^1 G_{\mathrm{rel}}^2\,\mathrm{d}z
=
\frac{\nu^2\,\kappa_T\,c^2}
{4\,g\,\alpha\,\beta\,\enthalpy\,d^6}
\int_0^1 \bigl[(DF_{\mathrm{rel}})^2
  + a^2 F_{\mathrm{rel}}^2\bigr]\,\mathrm{d}z,
\end{equation}
which, upon rearranging, yields
\begin{equation}
\label{eq:rel-2-186}
\boxed{
\Rarel
= \frac{g\,\alpha\,\beta\,d^4}{\kappa_T\,\nu}
  \cdot\frac{\enthalpy}{\rdensity\,c^2}
= \frac{\displaystyle\int_0^1
  \bigl[(DF_{\mathrm{rel}})^2 + a^2 F_{\mathrm{rel}}^2\bigr]
  \,\mathrm{d}z}
{a^2\;\displaystyle\int_0^1
  \frac{\enthalpy}{c^2}\,G_{\mathrm{rel}}^2\,\mathrm{d}z}.
}
\end{equation}
This is precisely the relativistic Rayleigh
quotient~\eqref{eq:rel-2-141}, confirming that the first variational
principle has the same physical content in the relativistic theory:

\emph{Relativistic thermal instability occurs at the minimum
temperature gradient at which a balance can be steadily maintained
between the kinetic energy dissipated by viscosity (acting on the
relativistic enthalpy-inertia $\enthalpy/c^2$) and the internal energy
released by the relativistic buoyancy force.}

\subsection{Relativistic entropy production and the BDNK
second law}\label{sec:rel-14d}

The BDNK theory provides a natural framework for connecting
the variational principle to entropy production.  In the BDNK
first-order formalism, the entropy current is
\begin{equation}
\label{eq:rel-entropy-current}
s^\mu = s\,n\,u^\mu + \frac{Q^\mu}{T}
  + \text{(frame-dependent first-order terms)},
\end{equation}
where $s$ is the entropy per baryon, $n$ the baryon number density,
and $Q^\mu$ is the BDNK heat flux including frame corrections.
Its divergence is (Bemfica et al.\ 2018; Kovtun 2019)
\begin{equation}
\label{eq:rel-entropy-production}
\nabla_\mu s^\mu
= -\frac{\bulkP^2}{\bulkvisc\,T}
  - \frac{\Pi_{\mu\nu}\Pi^{\mu\nu}}{2\,\shearvisc\,T}
  - \frac{Q_\mu Q^\mu}{\thermcond\,T^2}
  \;\geqslant\; 0.
\end{equation}
The second law $\nabla_\mu s^\mu \geqslant 0$ is guaranteed
provided the transport coefficients satisfy
$\shearvisc,\,\bulkvisc,\,\thermcond \geqslant 0$ and the BDNK
frame coefficients satisfy the coupled inequalities that ensure
strong hyperbolicity (Bemfica, Disconzi \& Noronha 2019).

At marginal stability ($\sigma = 0$), only the shear viscous and heat
flux contributions survive in the linearised perturbation.  The
volume-integrated entropy production rate in the perturbed layer
becomes
\begin{equation}
\label{eq:rel-entropy-rate}
\dot{S}_{\mathrm{prod}}
= \int \nabla_\mu s^\mu\,\sqrt{-g}\;\mathrm{d}^3 x
= \frac{1}{T}\bigl(\epsilon_v^{\mathrm{rel}}
  + \epsilon_{\mathrm{heat}}^{\mathrm{rel}}\bigr),
\end{equation}
where $\epsilon_{\mathrm{heat}}^{\mathrm{rel}}$ is the rate of entropy
production from heat conduction.  At the marginal state
$\epsilon_v^{\mathrm{rel}} = \epsilon_g^{\mathrm{rel}}$, and the
variational principle~\eqref{eq:rel-2-186} selects the perturbation
mode that \emph{maximises} the entropy production rate at the
threshold of instability, consistent with the principle of maximum
entropy production.

\begin{relcorrection}
The connection to entropy production is qualitatively the same as in
the classical theory but gains two important features in the
relativistic context:

\begin{enumerate}
\item In the BDNK first-order theory, the entropy current~\eqref{eq:rel-entropy-current}
  contains frame-dependent first-order gradient terms whose coefficients
  are chosen to ensure causal propagation and strong hyperbolicity.
  The entropy production inequality~\eqref{eq:rel-entropy-production}
  is maintained by the positivity of the transport coefficients
  together with the BDNK frame constraints.

\item The maximum entropy production principle, which classically
  determines the critical Rayleigh number, is promoted to a
  \emph{covariant} statement: the four-divergence of the entropy
  current is maximised among all perturbation modes satisfying the
  relativistic constraints.
\end{enumerate}
\end{relcorrection}

\begin{causalitycheck}
At marginal stability ($\sigma = 0$) all time derivatives vanish,
so the dissipative constitutive relations play no role in the onset
equations.  Consequently, the BDNK first-order theory and the
Israel--Stewart second-order theory yield \emph{identical} results
at the marginal state.  Away from marginality, the BDNK theory
ensures causal propagation through the coupled inequalities on the
transport and frame coefficients (Bemfica et al.\ 2018; Kovtun 2019),
which guarantee that all characteristic speeds of the BDNK system
satisfy
\begin{equation}
v_{\mathrm{char}}^2 \leqslant c^2.
\end{equation}
Unlike the Israel--Stewart theory, no relaxation times $\taupi$,
$\tauq$, $\tauPi$ are needed; causality is achieved through the
choice of hydrodynamic frame and constrained first-order transport
coefficients.
\end{causalitycheck}

\subsection{Non-relativistic recovery}\label{sec:rel-14e}

In the limit $v/c \to 0$ (equivalently
$\enthalpy = \edensity + p \to \rdensity c^2$), every relativistic
expression derived above reduces to its classical counterpart:

\begin{enumerate}
\item The enthalpy weight
  $\enthalpy/c^2 \to \rdensity$, so the relativistic Rayleigh
  quotient~\eqref{eq:rel-2-141} reduces to equation~(2-141).

\item The relativistic Rayleigh number
  $\Rarel = \Ra\cdot\enthalpy/(\rdensity c^2)
  \;\NRlimit\; \Ra$.

\item The orthogonality relation~\eqref{eq:rel-2-139} loses the
  weight $\enthalpy/c^2$ and becomes equation~(2-139).

\item The viscous dissipation rate~\eqref{eq:rel-2-178} reduces to
  $\epsilon_v = (\rho\nu/4d^2)\int_0^1 G^2\,\mathrm{d}z$,
  matching equation~(2-178).

\item The buoyancy energy release rate~\eqref{eq:rel-2-183} reduces
  to equation~(2-183).

\item The energy balance~\eqref{eq:rel-2-186} becomes
  $\Ra = \int[(DF)^2 + a^2 F^2]\,\mathrm{d}z\,/\,
  (a^2\int G^2\,\mathrm{d}z)$,
  reproducing equation~(2-186).

\item The BDNK entropy production
  inequality~\eqref{eq:rel-entropy-production} reduces to the
  classical Clausius--Duhem inequality in the non-relativistic limit,
  as the frame corrections vanish at $\order{v^2/c^2}$.
\end{enumerate}

\emph{All relativistic corrections are therefore of order
$\order{(\edensity + p)/(\rdensity c^2) - 1}
= \order{p/(\rdensity c^2)}$
and vanish identically in the non-relativistic limit.}
